\documentclass{cfc2023}
\usepackage{pslatex}
%\usepackage{amsmath}
%\usepackage{amsfonts}
%\usepackage{amssymb}
\thispagestyle{empty}

\title{Reduced Order Models on a Variational Multi-Scale Model of Navier--Stokes}

\author{Davide Torlo$^{*}$, Samuele Rubino$^{\dag}$ and Giovanni Stabile$^{\ddag}$}



\address{$^{*}$ SISSA mathLab, SISSA International School for Advanced Studies,\\
via Bonomea, 265, 34136 Trieste, Italy\\
e-mail:  davide.torlo@sissa.it, web page: https://www.davidetorlo.it
\and
$^{\dag}$ Departamento EDAN \& IMUS,\\ Universidad de Sevilla, Campus de Reina Mercedes,
41012, Sevilla, Spain \and 
$^{\ddag}$ SISSA mathLab, SISSA International School for Advanced Studies,\\
via Bonomea, 265, 34136 Trieste, Italy
}

\begin{document}
\thispagestyle{empty}

In this work, we study parametric incompressible flows given by Navier--Stokes equations. At the discrete level, we use continuous finite element (FE) method. 
The discrete model relies on a variational multi-scale (VMS) approach, which separates the large, sub-filter and small scales. 
The first two are resolved, while the last is not. 
The model includes sub-grid eddy viscosity to take care of the interaction between sub-filter and small scales.
A local projection stabilization (LPS) term is introduced onto the sub-filter terms to provide stability also in a FE context. 
This term is based on interpolation and projection operators that penalizes the oscillations on the sub-filter scale\cite{Rebollo15,Ahmed17}.\\
The computational costs of such simulations, though being faster than other models as the small scales are not resolved, are still large, in particular in the time dependent case.
Moreover, dealing with parametric problems in a multi-query context, where hundreds or thousands of simulations are necessary, makes the computational burden unbearable.
We propose a Galerkin projection of the equations onto a POD-generated reduced basis space as reduced order model (ROM).
To take care of the nonlinearities of the problem, different hyper-reduction techniques are studied in order to obtain a reduced model that is independent of the dimension of the FE space. \\
We provide simulations in two dimensions to validate the ROM and to prove the computational advantage of the approach.




\begin{center}
\begin{thebibliography}{99}
\bibitem{Rebollo15} T.~C.~Rebollo, M.~G.~Mármol and S. Rubino. Numerical analysis of a finite element projection-based VMS
turbulence model with wall laws. \emph{Comput. Methods Appl. Mech. Engrg.}, Vol. \textbf{285}, pp. 379--405, 2015.
\bibitem{Ahmed17} N.~Ahmed, T.~C.~Rebollo, V.~John and S. Rubino. Analysis of a Full Space--Time Discretization of the Navier--Stokes Equations by a
Local Projection Stabilization Method. \emph{IMA Journal of Numerical Analysis}, Vol. \textbf{37}, pp. 1437--1467, 2017.
\end{thebibliography}
\end{center}

\end{document}


